\clearpage
\phantomsection
\addcontentsline{toc}{section}{Referências}
\begin{thebibliography}{99}
\setlength{\itemsep}{3pt}

\bibitem{bruck1990} BRUCK, J. On the convergence properties of the Hopfield model.
\textbf{Proceedings of the IEEE}, v.~78, n.~10, p.~1579--1585, 1990.
DOI: \url{https://doi.org/10.1109/5.58341}.

\bibitem{bruckroychowdhury1990} BRUCK, J.; ROYCHOWDHURY, V. P. On the number of spurious memories in the
Hopfield model. \textbf{IEEE Transactions on Information Theory}, v.~36, n.~2, p.~393--397, 1990.
DOI: \url{https://doi.org/10.1109/18.52486}.

\bibitem{chandola2009} CHANDOLA, V.; BANERJEE, A.; KUMAR, V. Anomaly detection: a survey.
\textbf{ACM Computing Surveys}, v.~41, n.~3, art.~15, 2009.
DOI: \url{https://doi.org/10.1145/1541880.1541882}.

\bibitem{demircigil2017} DEMIRCIGIL, M. \textit{et al.} On a model of associative memory with huge storage
capacity. \textbf{Journal of Statistical Physics}, v.~168, n.~2, p.~288--299, 2017.
DOI: \url{https://doi.org/10.1007/s10955-017-1806-y}.

\bibitem{fawcett2006} FAWCETT, T. An introduction to ROC analysis. \textbf{Pattern Recognition Letters},
v.~27, n.~8, p.~861--874, 2006. DOI: \url{https://doi.org/10.1016/j.patrec.2005.10.010}.

\bibitem{hertz1991} HERTZ, J.; KROGH, A.; PALMER, R. G. \textbf{Introduction to the Theory of Neural
Computation}. Redwood City: Addison-Wesley, 1991.

\bibitem{hopfield1982} HOPFIELD, J. J. Neural networks and physical systems with emergent collective
computational abilities. \textbf{Proceedings of the National Academy of Sciences (PNAS)}, v.~79, n.~8,
p.~2554--2558, 1982. DOI: \url{https://doi.org/10.1073/pnas.79.8.2554}.

\bibitem{hopfield1984} HOPFIELD, J. J. Neurons with graded response have collective computational properties
like those of two-state neurons. \textbf{Proceedings of the National Academy of Sciences (PNAS)}, v.~81,
n.~10, p.~3088--3092, 1984. DOI: \url{https://doi.org/10.1073/pnas.81.10.3088}.

\bibitem{hopfieldtank1985} HOPFIELD, J. J.; TANK, D. W. ``Neural'' computation of decisions in optimization
problems. \textbf{Biological Cybernetics}, v.~52, n.~3, p.~141--152, 1985.
DOI: \url{https://doi.org/10.1007/BF00339943}.

\bibitem{kim2022} KIM, S.; CHOI, K.; CHOI, H.-S.; LEE, B.; YOON, S. Towards a rigorous evaluation of
time-series anomaly detection. \textbf{Proceedings of the AAAI Conference on Artificial Intelligence},
v.~36, n.~7, p.~7194--7201, 2022. DOI: \url{https://doi.org/10.1609/aaai.v36i7.20680}.

\bibitem{krotov2016} KROTOV, D.; HOPFIELD, J. J. Dense associative memory for pattern recognition. In:
\textbf{Advances in Neural Information Processing Systems 29 (NeurIPS 2016)}, 2016, p.~1172--1180.
arXiv:1606.01164.

\bibitem{mceliece1987} MCELIECE, R. J.; POSNER, E. C.; RODEMICH, E. R.; VENKATESH, S. S. The capacity of the
Hopfield associative memory. \textbf{IEEE Transactions on Information Theory}, v.~33, n.~4, p.~461--482,
1987. DOI: \url{https://doi.org/10.1109/TIT.1987.1057328}.

\bibitem{hopfieldlayers} ML-JKU. \textbf{hopfield-layers}: código-fonte das camadas Hopfield, indicado por
Ramsauer \textit{et al.} (2021). Disponível em: \url{https://github.com/ml-jku/hopfield-layers}.

\bibitem{omnianomaly} NETMANAIOPS. \textbf{OmniAnomaly}: código e \textit{Server Machine Dataset} (SMD).
Repositório do GitHub (licença MIT declarada no repositório). Disponível em:
\url{https://github.com/NetManAIOps/OmniAnomaly}. Acesso em: 6 nov. 2026.

\bibitem{pedregosa2011} PEDREGOSA, F. \textit{et al.} Scikit-learn: machine learning in Python.
\textbf{Journal of Machine Learning Research}, v.~12, p.~2825--2830, 2011.

\bibitem{ramsauer2021} RAMSAUER, H. \textit{et al.} Hopfield networks is all you need. In:
\textbf{International Conference on Learning Representations (ICLR)}, 2021. arXiv:2008.02217.

\bibitem{saito2015} SAITO, T.; REHMSMEIER, M. The precision-recall plot is more informative than the ROC plot
when evaluating binary classifiers on imbalanced datasets. \textbf{PLOS ONE}, v.~10, n.~3, e0118432, 2015.
DOI: \url{https://doi.org/10.1371/journal.pone.0118432}.

\bibitem{storkey1997} STORKEY, A. J. Increasing the capacity of a Hopfield network without sacrificing
functionality. In: \textbf{Artificial Neural Networks --- ICANN'97}. Lecture Notes in Computer Science.
Berlin: Springer, 1997, p.~451--456. DOI: \url{https://doi.org/10.1007/BFb0020196}.

\bibitem{su2019} SU, Y. \textit{et al.} Robust anomaly detection for multivariate time series through
stochastic recurrent neural network. In: \textbf{Proceedings of the 25th ACM SIGKDD International Conference
on Knowledge Discovery \& Data Mining (KDD '19)}. ACM, 2019.
DOI: \url{https://doi.org/10.1145/3292500.3330672}.

\bibitem{enunciado} TRABALHO de Implementação de Redes Neurais Artificiais: descrição do trabalho. Documento
da disciplina fornecido pelo professor (PDF), [s.d.].

\bibitem{wikihopfield} WIKIPEDIA. \textbf{Hopfield network}. Material fornecido pelo professor (impressão de
14 out. 2026; versão \texttt{oldid=1377477270}). Disponível em:
\url{https://en.wikipedia.org/wiki/Hopfield_network}.

\end{thebibliography}
